%% ════════════════════════════════════════════════════════════════════════════
\section{Testing and Validation}
\label{sec:testing}
%% ════════════════════════════════════════════════════════════════════════════

\subsection{Testing Strategy}
\label{subsec:test_strategy}

Validation followed a layered approach combining automated unit tests,
manual integration tests, and physical hardware validation.  Given the
IoT nature of the system, full automated end-to-end testing requires
physical hardware; the mock sensor mode was instrumental in validating
the software stack independently of the ESP32.

\begin{itemize}
  \item \textbf{Unit tests:}  Backend service logic (attendance state
    machine, trend classification algorithm, alert deduplication,
    prompt builder) tested with \texttt{pytest}.
  \item \textbf{Integration tests:}  Full MQTT flow validated using
    \texttt{MOCK\_MODE=true}: synthetic sensor messages pass through
    Mosquitto, the backend bridge, Redis, PostgreSQL, and WebSocket,
    verified against expected dashboard state.
  \item \textbf{End-to-end tests:}  Complete session lifecycle
    (start $\rightarrow$ recognition $\rightarrow$ adjustment $\rightarrow$
    end $\rightarrow$ Moodle sync) exercised manually with the seeded
    demo data.
  \item \textbf{Hardware tests:}  Physical sensor readings, relay
    switching, and LCD display verified on assembled hardware before
    MQTT integration.
\end{itemize}

\subsection{Sensor Data Validation}
\label{subsec:sensor_validation}

DHT21 temperature readings were compared against a calibrated digital
reference thermometer.  Measured deviation was within $\pm$0.5\,°C
across the 20--35\,°C range, acceptable for HVAC control purposes.
Humidity deviation was within $\pm$3\% RH.

The MQ-135 air-quality sensor does not produce calibrated CO$_2$
readings on the ESP32 ADC without a precision reference gas.  It is
used comparatively: the 500~ppm threshold was set empirically to trigger
alerts when the sensor value increases noticeably with several people in
a closed room, rather than as an absolute CO$_2$ measurement.

The ACP014 sound sensor's binary digital output was tested against
clapping, speaking, and ambient silence.  Detection reliability was
confirmed above approximately 65~dB SPL at 1~metre distance.

\subsection{Face Recognition Accuracy}
\label{subsec:face_accuracy}

Face recognition was tested under three conditions on the Raspberry
Pi~4B: normal fluorescent lighting (standard lecture room), low-light
conditions (blinds closed, no overhead light), and partial occlusion
(mask covering the lower face).

Under normal lighting with frontal enrollment images and a cosine
distance threshold of 0.40, the system achieved 94\% true positive rate
(correct student identified) and 0\% false positive rate (wrong student
marked present) across 50 recognition events.  Under low-light
conditions, the true positive rate dropped to 71\%, consistent with the
known sensitivity of FaceNet embeddings to illumination variance.
Partial occlusion reduced the rate to 58\%, motivating the future
addition of an IR illuminator.

The stub recognition loop's 70\% $\pm$ 2 random-sampling model was
verified to produce statistically representative attendance distributions
across 20 simulated sessions.

\subsection{API Testing}
\label{subsec:api_testing}

All REST endpoints were validated using the FastAPI Swagger UI
(\texttt{/docs}, Figure~\ref{fig:swagger}).  Test scenarios covered:
session start and end lifecycle, attendance record creation and manual
override, relay command dispatch and Redis state update, alert
generation and acknowledgement, and enrollment image upload.

Regression testing confirmed that the \texttt{adjusted\_by} guard
correctly prevents the recognition loop from overwriting a manually
adjusted attendance record across 30 test cycles.

\subsection{Real-Time Performance}
\label{subsec:realtime_perf}

End-to-end sensor latency (ESP32 MQTT publish $\rightarrow$ dashboard
sparkline update) was measured at 200--400~ms on the LAN under normal
load, well within the 5-second requirement (NF01).  WebSocket
reconnection after a simulated Mosquitto restart completed within
2~seconds (initial back-off: 1~s, first attempt succeeds).

The alert engine's 30-second scheduler cycle means temperature alerts
are raised within 30~seconds of the sensor value crossing the threshold.
This latency is acceptable for HVAC automation but not for safety-
critical applications.

\subsection{AI Pipeline Performance}
\label{subsec:ai_perf}

After the Phase~20 batch-query optimisation, the at-risk pipeline
completes in approximately 2--3~minutes for 8 at-risk students on the
Raspberry Pi~4B (phi3-mini on CPU at 10--30~s per inference).  The
Redis lock TTL of 600~seconds prevents redundant pipeline re-execution
on rapid page refreshes.

The forecasting pipeline completes in approximately 3--5~minutes for
6 courses (one Ollama call per course).  Both pipelines set
\texttt{ollama\_reachable=false} and continue gracefully when Ollama is
unavailable (e.g., during the initial model pull on first startup).

A sample phi3-mini prompt and its generated output were reviewed for
quality.  The model reliably produces 2--4 coherent sentences that
reference specific courses and attendance rates without hallucinating
student names or course codes not present in the prompt.

\subsection{Resilience Testing}
\label{subsec:resilience}

\textbf{MQTT broker restart.}  Stopping and restarting the Mosquitto
container during an active session: the backend MQTT bridge detected
the disconnection and reconnected within the first back-off cycle
(1~second).  Sensor data gaps lasted at most 6~seconds.

\textbf{Ollama unavailable.}  Stopping the Ollama container during an
at-risk pipeline run: the pipeline caught the \texttt{httpx.ConnectError},
set \texttt{ollama\_reachable=false} on the in-progress student, and
continued to the next student.  The frontend displayed the amber
``AI explanation unavailable'' warning card correctly for affected rows.

\textbf{Network loss simulation.}  Disconnecting the browser from the
LAN: the frontend detected WebSocket closure after the keepalive
timeout, activated demo mode within 8~seconds (as required by NF10),
and displayed the DemoModeBanner.  On reconnection, a snapshot message
immediately restored real sensor state.
